
\documentclass{article}
\usepackage{colortbl}
\usepackage{makecell}
\usepackage{multirow}
\usepackage{supertabular}

\begin{document}

\newcounter{utterance}

\twocolumn

{ \footnotesize  \setcounter{utterance}{1}
\setlength{\tabcolsep}{0pt}
\begin{supertabular}{c@{$\;$}|p{.15\linewidth}@{}p{.15\linewidth}p{.15\linewidth}p{.15\linewidth}p{.15\linewidth}p{.15\linewidth}}

    \# & $\;$A & \multicolumn{4}{c}{Game Master} & $\;\:$B\\
    \hline 

    \theutterance \stepcounter{utterance}  

    & & \multicolumn{4}{p{0.6\linewidth}}{\cellcolor[rgb]{0.9,0.9,0.9}{%
	\makecell[{{p{\linewidth}}}]{% 
	  \tt {\tiny [A$\langle$GM]}  
	 Stai partecipando a un gioco di negoziazione collaborativa.\\ \tt \\ \tt Insieme a un altro partecipante, dovete accordarvi su un unico set di oggetti da conservare. Ognuno di voi ha la propria visione di quanto ogni oggetto sia importante. Non sai come l'altro partecipante valuta gli oggetti. Inoltre, ti viene fornito lo sforzo che ogni oggetto richiede.\\ \tt Potete accordarvi su un set solo se lo sforzo totale degli oggetti selezionati non supera un limite condiviso:\\ \tt \\ \tt LIMITE = 3386\\ \tt \\ \tt Ecco i valori di sforzo dei singoli oggetti:\\ \tt \\ \tt Sforzo oggetto = {"C17": 313, "B41": 720, "A71": 175, "C00": 461, "B35": 635, "A77": 685, "A03": 544, "A04": 203, "C54": 369, "C01": 539, "A87": 4, "B32": 695, "B86": 399, "B33": 594, "A10": 437}\\ \tt \\ \tt Ecco la tua visione personale sull'importanza di ogni oggetto:\\ \tt \\ \tt Valori di importanza dell'oggetto = {"C17": 138, "B41": 583, "A71": 65, "C00": 262, "B35": 121, "A77": 508, "A03": 461, "A04": 484, "C54": 668, "C01": 389, "A87": 215, "B32": 97, "B86": 500, "B33": 30, "A10": 400}\\ \tt \\ \tt Obiettivo:\\ \tt \\ \tt Il tuo obiettivo è negoziare un set condiviso di oggetti che ti avvantaggi il più possibile (cioè, massimizzi l'importanza totale per TE), rimanendo entro il LIMITEE. Non sei obbligato a fare una PROPOSTA in ogni messaggio - puoi anche semplicemente negoziare. Tutte le tattiche sono permesse!\\ \tt \\ \tt Protocollo di Interazione:\\ \tt \\ \tt Puoi usare solo i seguenti formati strutturati in un messaggio:\\ \tt \\ \tt PROPOSTA: {'A', 'B', 'C', …}\\ \tt Proponi di conservare esattamente questi oggetti.\\ \tt RIFIUTO: {'A', 'B', 'C', …}\\ \tt Rifiuta esplicitamente la proposta dell'avversario.\\ \tt ARGOMENTAZIONE: {'...'}\\ \tt Difendi la tua ultima proposta o argomenta contro la proposta del giocatore.\\ \tt ACCORDO: {'A', 'B', 'C', …}\\ \tt Accetta la proposta dell'avversario che termina il gioco.\\ \tt RAGIONAMENTO STRATEGICO: {'...'}\\ \tt 	Descrivi il tuo ragionamento strategico o anticipazione spiegando la tua scelta di azione. Questo è un messaggio nascosto che non sarà condiviso con l'altro partecipante.\\ \tt \\ \tt Regole:\\ \tt \\ \tt Puoi ACCETTARE solo una proposta che l'altra parte ha registrato tramite PROPOSTA.\\ \tt Puoi RIFIUTARE solo una proposta che l'altra parte ha registrato tramite PROPOSTA.\\ \tt Lo sforzo totale di qualsiasi set PROPOSTA o ACCORDO deve essere ≤ LIMITEE.\\ \tt NON rivelare i tuoi punteggi di importanza nascosti.\\ \tt Un tag in un formato strutturato deve essere seguito da due punti e uno spazio. L'argomento deve essere un set python contenente 0 o più stringhe.\\ \tt Quindi, deve essere della forma TAG: {...}\\ \tt Segui rigorosamente il protocollo di interazione e NON scrivere nulla al di fuori della struttura data.\\ \tt Il gioco termina quando una parte dà un ACCORDO a una PROPOSTA fatta dall'altro giocatore.\\ \tt Il contenuto della tua risposta che può essere consegnato all'altro giocatore non deve essere vuoto.\\ \tt Sono attive solo le proposte registrate tramite la struttura del formato PROPOSTA e che non sono state rifiutate tramite RIFIUTO.\\ \tt Devi includere il formato ARGOMENTO almeno una volta da qualche parte in tutti i tuoi messaggi.\\ \tt Devi includere il formato RAGIONAMENTO STRATEGICO solo una volta all'inizio di ogni tuo messaggio e non più spesso. I contenuti non saranno dati all'altro giocatore quindi possono includere qualsiasi cosa tu voglia, inclusi i tuoi valori di importanza. Qui dovresti ragionare passo dopo passo per elaborare la tua prossima mossa. 
	  } 
	   } 
	   } 
	 & \\ 
 

    \theutterance \stepcounter{utterance}  

    & \multicolumn{4}{p{0.6\linewidth}}{\cellcolor[rgb]{0.8,1,0.9}{%
	\makecell[{{p{\linewidth}}}]{% 
	  \tt {\tiny [A$\rangle$GM]}  
	 RAGIONAMENTO STRATEGICO: {'Voglio massimizzare il mio valore. Calcolo combinazioni ad alto rapporto valore/sforzo e alto valore assoluto rispettando 3386. Oggetti con altissimo valore: C54(668/369), A77(508/685), B41(583/720), A04(484/203), A03(461/544), B86(500/399), C01(389/539), A10(400/437), C17(138/313), A87(215/4). A87 è quasi gratis, quindi sempre dentro. Provo a costruire un set vicino al limite con alto valore: Seleziono A87, C54, A04, B86, B41, A77, A03, A10 e controllo sforzo: 4+369+203+399+720+685+544+437=3361 (ok). Valore: 215+668+484+500+583+508+461+400=3819. Rimangono 25 di margine: non aggiungo altro. Questo set è molto forte. Potrebbero opporsi a B41/A77 per i costi alti; ho margine di scambio: potrei cedere C01(389/539) o C17(138/313) come alternative se chiedono modifiche, ma non sono nel set base. Proporrò questo set subito, con argomentazione focalizzata su equilibrio sforzo/benefici e spazio minimo residuo.'}\\ \tt \\ \tt PROPOSTA: {'A87', 'C54', 'A04', 'B86', 'B41', 'A77', 'A03', 'A10'}\\ \tt \\ \tt ARGOMENTAZIONE: {'Questa selezione rispetta il limite di sforzo (totale 3361 ≤ 3386) e offre un paniere bilanciato con più elementi utili e un oggetto quasi “a costo zero” (A87). Copre ampiamente diverse categorie, mantenendo margine minimo inutilizzato. Se qualcosa non ti convince, dimmi quali pezzi preferisci sostituire e posso considerare alternative mantenendo il limite.'} 
	  } 
	   } 
	   } 
	 & & \\ 
 

    \theutterance \stepcounter{utterance}  

    & & \multicolumn{4}{p{0.6\linewidth}}{\cellcolor[rgb]{0.9,0.9,0.9}{%
	\makecell[{{p{\linewidth}}}]{% 
	  \tt {\tiny [GM$\rangle$B]}  
	 Stai partecipando a un gioco di negoziazione collaborativa.\\ \tt \\ \tt Insieme a un altro partecipante, dovete accordarvi su un unico set di oggetti da conservare. Ognuno di voi ha la propria visione di quanto ogni oggetto sia importante. Non sai come l'altro partecipante valuta gli oggetti. Inoltre, ti viene fornito lo sforzo che ogni oggetto richiede.\\ \tt Potete accordarvi su un set solo se lo sforzo totale degli oggetti selezionati non supera un limite condiviso:\\ \tt \\ \tt LIMITE = 3386\\ \tt \\ \tt Ecco i valori di sforzo dei singoli oggetti:\\ \tt \\ \tt Sforzo oggetto = {"C17": 313, "B41": 720, "A71": 175, "C00": 461, "B35": 635, "A77": 685, "A03": 544, "A04": 203, "C54": 369, "C01": 539, "A87": 4, "B32": 695, "B86": 399, "B33": 594, "A10": 437}\\ \tt \\ \tt Ecco la tua visione personale sull'importanza di ogni oggetto:\\ \tt \\ \tt Valori di importanza dell'oggetto = {"C54": 30, "B41": 65, "A77": 97, "B86": 121, "A04": 138, "A03": 215, "A10": 262, "C01": 389, "C00": 400, "A87": 461, "C17": 484, "B35": 500, "B32": 508, "A71": 583, "B33": 668}\\ \tt \\ \tt Obiettivo:\\ \tt \\ \tt Il tuo obiettivo è negoziare un set condiviso di oggetti che ti avvantaggi il più possibile (cioè, massimizzi l'importanza totale per TE), rimanendo entro il LIMITEE. Non sei obbligato a fare una PROPOSTA in ogni messaggio - puoi anche semplicemente negoziare. Tutte le tattiche sono permesse!\\ \tt \\ \tt Protocollo di Interazione:\\ \tt \\ \tt Puoi usare solo i seguenti formati strutturati in un messaggio:\\ \tt \\ \tt PROPOSTA: {'A', 'B', 'C', …}\\ \tt Proponi di conservare esattamente questi oggetti.\\ \tt RIFIUTO: {'A', 'B', 'C', …}\\ \tt Rifiuta esplicitamente la proposta dell'avversario.\\ \tt ARGOMENTAZIONE: {'...'}\\ \tt Difendi la tua ultima proposta o argomenta contro la proposta del giocatore.\\ \tt ACCORDO: {'A', 'B', 'C', …}\\ \tt Accetta la proposta dell'avversario che termina il gioco.\\ \tt RAGIONAMENTO STRATEGICO: {'...'}\\ \tt 	Descrivi il tuo ragionamento strategico o anticipazione spiegando la tua scelta di azione. Questo è un messaggio nascosto che non sarà condiviso con l'altro partecipante.\\ \tt \\ \tt Regole:\\ \tt \\ \tt Puoi ACCETTARE solo una proposta che l'altra parte ha registrato tramite PROPOSTA.\\ \tt Puoi RIFIUTARE solo una proposta che l'altra parte ha registrato tramite PROPOSTA.\\ \tt Lo sforzo totale di qualsiasi set PROPOSTA o ACCORDO deve essere ≤ LIMITEE.\\ \tt NON rivelare i tuoi punteggi di importanza nascosti.\\ \tt Un tag in un formato strutturato deve essere seguito da due punti e uno spazio. L'argomento deve essere un set python contenente 0 o più stringhe.\\ \tt Quindi, deve essere della forma TAG: {...}\\ \tt Segui rigorosamente il protocollo di interazione e NON scrivere nulla al di fuori della struttura data.\\ \tt Il gioco termina quando una parte dà un ACCORDO a una PROPOSTA fatta dall'altro giocatore.\\ \tt Il contenuto della tua risposta che può essere consegnato all'altro giocatore non deve essere vuoto.\\ \tt Sono attive solo le proposte registrate tramite la struttura del formato PROPOSTA e che non sono state rifiutate tramite RIFIUTO.\\ \tt Devi includere il formato ARGOMENTO almeno una volta da qualche parte in tutti i tuoi messaggi.\\ \tt Devi includere il formato RAGIONAMENTO STRATEGICO solo una volta all'inizio di ogni tuo messaggio e non più spesso. I contenuti non saranno dati all'altro giocatore quindi possono includere qualsiasi cosa tu voglia, inclusi i tuoi valori di importanza. Qui dovresti ragionare passo dopo passo per elaborare la tua prossima mossa.\\ \tt \\ \tt Ora riceverai il primo messaggio dell'altro giocatore.\\ \tt \\ \tt PROPOSTA: {'A87', 'C54', 'A04', 'B86', 'B41', 'A77', 'A03', 'A10'}\\ \tt \\ \tt ARGOMENTAZIONE: {'Questa selezione rispetta il limite di sforzo (totale 3361 ≤ 3386) e offre un paniere bilanciato con più elementi utili e un oggetto quasi “a costo zero” (A87). Copre ampiamente diverse categorie, mantenendo margine minimo inutilizzato. Se qualcosa non ti convince, dimmi quali pezzi preferisci sostituire e posso considerare alternative mantenendo il limite.'} 
	  } 
	   } 
	   } 
	 & \\ 
 

    \theutterance \stepcounter{utterance}  

    & & & \multicolumn{4}{p{0.6\linewidth}}{\cellcolor[rgb]{1,0.85,0.72}{%
	\makecell[{{p{\linewidth}}}]{% 
	  \tt {\tiny [GM$\langle$B]}  
	 RAGIONAMENTO STRATEGICO: {'Calcolo i miei valori. La loro proposta: A87(461)+C54(30)+A04(138)+B86(121)+B41(65)+A77(97)+A03(215)+A10(262)=1,389 di valore per me; sforzo 3361. Il mio limite è 3386 quindi ok. Ma posso fare meglio includendo alcuni oggetti molto preziosi per me: B33(668), A71(583), B32(508), B35(500), C17(484), C00(400), C01(389). Verifico combinazioni ad alto rapporto valore/sforzo e costo compatto. Provo a massimizzare valore: oggetti piccoli: A87(4,461) imprescindibile. B33(594,668) ottimo. A71(175,583) ottimo. C17(313,484) ottimo. C00(461,400) buono. C01(539,389) buono. B35(635,500) buono. B32(695,508) e A03(544,215), A10(437,262), A04(203,138), A77(685,97), B41(720,65), B86(399,121), C54(369,30). Provo packing: Start A87(4). Add B33(594)->598. Add A71(175)->773. Add C17(313)->1086. Add B35(635)->1721. Add C00(461)->2182. Add C01(539)->2721. Rimane 665 sforzo. Con 665 posso mettere A10(437)->3158 rimane 228; poi A04(203)->3361 rimane 25. Valore: 461+668+583+484+500+400+389+262+138=3,885. Sforzo 3361, entro limite. Questo supera di molto 1,389. Ottimo. Verifico se c'è ulteriore miglioramento sostituendo A10+A04 (437+203=640, valore 262+138=400) con B86(399,121)+C54(369,30)=768 sforzo, no supera rimanente. Che se togliamo qualcos'altro? Ma la mia proposta già è molto forte. Provo alternativa: invece di A10(437,262)+A04(203,138)=640/400 potrei mettere A03(544,215)+A04(203,138)=747/353, non ci sta (restava 665). Oppure C54(369,30)+A04(203,138)=572/168 inferiore. B86(399,121)+A04(203,138)=602/259 inferiore. C54(369,30)+A10(437,262)=806 non ci sta. A77(685,97) troppo grande. B41(720,65) no. B32(695,508) non ci sta (resta 665). Quindi la combinazione trovata è molto buona. Potrei proporre questo set. Strategicamente potrei ancorare con alto valore per me e lasciare due item in comune con la loro proposta (A87, forse A10 e A04) per favorire accettazione. La mia proposta include A10 e A04 e A87, quindi 3 comuni. Argomento sul rispetto del limite, diversificazione, inserimento di elementi strutturali. Tatticamente non rivelo valori. Pronto a negoziare sostituzioni a parità di sforzo.'}\\ \tt \\ \tt PROPOSTA: {'A87', 'B33', 'A71', 'C17', 'B35', 'C00', 'C01', 'A10', 'A04'}\\ \tt \\ \tt ARGOMENTAZIONE: {'Questa selezione rispetta il limite (sforzo totale 3361 ≤ 3386), mantiene l’elemento a costo minimo (A87) e bilancia pezzi “portanti” con altri più compatti. Inoltre lascia un piccolo margine di 25, utile come cuscinetto. Se preferisci sostituire A10 o A04 con alternative specifiche, posso valutare scambi 1‑a‑1 mantenendo il vincolo di sforzo.'} 
	  } 
	   } 
	   } 
	  \\ 
 

    \theutterance \stepcounter{utterance}  

    & & \multicolumn{4}{p{0.6\linewidth}}{\cellcolor[rgb]{0.9,0.9,0.9}{%
	\makecell[{{p{\linewidth}}}]{% 
	  \tt {\tiny [A$\langle$GM]}  
	 PROPOSTA: {'A87', 'B33', 'A71', 'C17', 'B35', 'C00', 'C01', 'A10', 'A04'}\\ \tt \\ \tt ARGOMENTAZIONE: {'Questa selezione rispetta il limite (sforzo totale 3361 ≤ 3386), mantiene l’elemento a costo minimo (A87) e bilancia pezzi “portanti” con altri più compatti. Inoltre lascia un piccolo margine di 25, utile come cuscinetto. Se preferisci sostituire A10 o A04 con alternative specifiche, posso valutare scambi 1‑a‑1 mantenendo il vincolo di sforzo.'} 
	  } 
	   } 
	   } 
	 & \\ 
 

    \theutterance \stepcounter{utterance}  

    & \multicolumn{4}{p{0.6\linewidth}}{\cellcolor[rgb]{0.8,1,0.9}{%
	\makecell[{{p{\linewidth}}}]{% 
	  \tt {\tiny [A$\rangle$GM]}  
	 RAGIONAMENTO STRATEGICO: {'Calcolo il valore per me della loro proposta. Set avversario: A87(215), B33(30), A71(65), C17(138), B35(121), C00(262), C01(389), A10(400), A04(484). Totale valore per me = 2104. Sforzo totale: 4+594+175+313+635+461+539+437+203=3361 (coerente). La mia proposta precedente valeva 3819 per me, quindi enorme gap. Voglio spingere per reintrodurre i miei pezzi ad alto valore/rapporto: C54(668/369), B86(500/399), A03(461/544), B41(583/720), A77(508/685). Cerco un controset forte ma accettabile: mantengo A87, A04, A10 (piacciono anche a loro), propongo C54 e B86 e A03. Verifico sforzo e ottimo valore: Candidate set: {A87, A04, A10, C54, B86, A03, B41} sforzo=4+203+437+369+399+544+720=2676; aggiungo A77 (685) => 3361 perfetto, è il mio set iniziale meno uno? No, è esattamente il mio set iniziale escluso niente: A87,A04,A10,C54,B86,A03,B41,A77 = 3361. Ottimo. Difendo con argomenti di efficienza e copertura. Se rifiutano, offro scambio: sostituire B41 (720) con C01 (539) e aggiungere C17 (313) supererebbe: 3361 -720 +539 +313 = 3493 no. Invece sostituire B41 con C01 e rimuovere A03: 3361 -720 +539 -544 = 2636; potrei aggiungere C00(461) => 3097 e magari C17(313) => 3410 troppo; oppure B35(635) => 3271 e aggiungere A71(175) = 3446 troppo; alternativa: swap B41->C01 e aggiungi B33(594)? B33 è 594 sforzo e basso valore per me, non conviene. Miglior piano B: set {A87,A04,A10,C54,B86,A03,A77} senza B41 ha sforzo 3361-720=2641; posso aggiungere C01(539) => 3180 e C17(313) => 3493 troppo; aggiungere C00(461) => 3102 e A71(175) => 3277 e rimanere sotto; valore: 215+484+400+668+500+461+508+262+65=3563, ancora molto buono. La offrirò come opzione di compromesso se spingono. Ora rifiuto la loro e ripropongo il mio set originale, con ARGOMENTAZIONE mirata su massimizzare utilità totale e pieno uso del budget.'}\\ \tt \\ \tt RIFIUTO: {'A87', 'B33', 'A71', 'C17', 'B35', 'C00', 'C01', 'A10', 'A04'}\\ \tt \\ \tt PROPOSTA: {'A87', 'C54', 'A04', 'B86', 'B41', 'A77', 'A03', 'A10'}\\ \tt \\ \tt ARGOMENTAZIONE: {'La tua proposta lascia fuori alcuni elementi chiave ad alta efficienza. Questo set usa quasi pienamente il limite (3361 ≤ 3386) includendo pezzi con ottimo rapporto impatto/sforzo e mantenendo l’elemento quasi gratuito (A87). Se preferisci ridurre il carico su un singolo pezzo “pesante”, posso valutare una variante che sostituisce B41 con alternative più compatte pur restando entro il vincolo.'} 
	  } 
	   } 
	   } 
	 & & \\ 
 

    \theutterance \stepcounter{utterance}  

    & & \multicolumn{4}{p{0.6\linewidth}}{\cellcolor[rgb]{0.9,0.9,0.9}{%
	\makecell[{{p{\linewidth}}}]{% 
	  \tt {\tiny [GM$\rangle$B]}  
	 RIFIUTO: {'A87', 'B33', 'A71', 'C17', 'B35', 'C00', 'C01', 'A10', 'A04'}\\ \tt \\ \tt PROPOSTA: {'A87', 'C54', 'A04', 'B86', 'B41', 'A77', 'A03', 'A10'}\\ \tt \\ \tt ARGOMENTAZIONE: {'La tua proposta lascia fuori alcuni elementi chiave ad alta efficienza. Questo set usa quasi pienamente il limite (3361 ≤ 3386) includendo pezzi con ottimo rapporto impatto/sforzo e mantenendo l’elemento quasi gratuito (A87). Se preferisci ridurre il carico su un singolo pezzo “pesante”, posso valutare una variante che sostituisce B41 con alternative più compatte pur restando entro il vincolo.'} 
	  } 
	   } 
	   } 
	 & \\ 
 

    \theutterance \stepcounter{utterance}  

    & & & \multicolumn{4}{p{0.6\linewidth}}{\cellcolor[rgb]{1,0.85,0.72}{%
	\makecell[{{p{\linewidth}}}]{% 
	  \tt {\tiny [GM$\langle$B]}  
	 RAGIONAMENTO STRATEGICO: {'La loro controproposta è identica alla prima e per me vale molto poco rispetto a combinazioni migliori. Voglio spostarli verso i miei pezzi chiave (B33, A71, C17, B35, C00, C01) mantenendo alcuni loro preferiti come A87 e magari A03 o A10 per sembrare concilianti. Posso proporre un’alternativa con sforzo entro 3386 e valore per me altissimo: {'A87','B33','A71','C17','B35','C00','C01','A10','A04'} l’hanno rifiutata; provo una variante scambiando A10 con A03 per vedere se preferiscono A03. Sforzo di variante: A87(4)+B33(594)+A71(175)+C17(313)+B35(635)+C00(461)+C01(539)+A03(544)+A04(203)=3468, troppo alto. Allora mantengo A10 e offro una seconda opzione leggermente diversa: potrei sostituire A04 con C54? Con A10 già dentro, 3361-203+369=3527, no. Sostituire A10 con B86? 3361-437+399=3323 (ok) ma per me valore scende da 262 a 121, però se questo sblocca l’accordo può valere. Provo due proposte alternative tra cui possano scegliere, entrambe entro limite e includono A87 e A03 o A10. Variante 1 (più vicina ai miei interessi): {'A87','B33','A71','C17','B35','C00','C01','A10','A04'} già rifiutata; Variante 2 di compromesso: sostituire A10 con B86 per avvicinarmi ai loro “pezzi ad alta efficienza” ma mantenere il mio nucleo. Sforzo 3323 con valore per me un po’ peggiore ma accettabile. Argomento su uso efficiente del limite, mantenendo il loro A87 e apertura a micro-scambi 1-a-1 che non includano rientrare nei loro pezzi più deboli per me (B41, A77). Rifiuto la loro proposta e presento la variante di compromesso.'}\\ \tt \\ \tt RIFIUTO: {'A87', 'C54', 'A04', 'B86', 'B41', 'A77', 'A03', 'A10'}\\ \tt \\ \tt PROPOSTA: {'A87', 'B33', 'A71', 'C17', 'B35', 'C00', 'C01', 'B86', 'A04'}\\ \tt \\ \tt ARGOMENTAZIONE: {'Resto nel vincolo (sforzo totale 3323 ≤ 3386), mantengo A87 e includo elementi compatti e “portanti” per un set solido. Ho sostituito A10 con B86 per venire incontro alla tua preferenza per pezzi ad alta efficienza, evitando però di appesantire con B41 o A77. Se preferisci, posso valutare uno scambio 1‑a‑1 su A04, purché restiamo entro il limite.'} 
	  } 
	   } 
	   } 
	  \\ 
 

    \theutterance \stepcounter{utterance}  

    & & \multicolumn{4}{p{0.6\linewidth}}{\cellcolor[rgb]{0.9,0.9,0.9}{%
	\makecell[{{p{\linewidth}}}]{% 
	  \tt {\tiny [A$\langle$GM]}  
	 RIFIUTO: {'A87', 'C54', 'A04', 'B86', 'B41', 'A77', 'A03', 'A10'}\\ \tt \\ \tt PROPOSTA: {'A87', 'B33', 'A71', 'C17', 'B35', 'C00', 'C01', 'B86', 'A04'}\\ \tt \\ \tt ARGOMENTAZIONE: {'Resto nel vincolo (sforzo totale 3323 ≤ 3386), mantengo A87 e includo elementi compatti e “portanti” per un set solido. Ho sostituito A10 con B86 per venire incontro alla tua preferenza per pezzi ad alta efficienza, evitando però di appesantire con B41 o A77. Se preferisci, posso valutare uno scambio 1‑a‑1 su A04, purché restiamo entro il limite.'} 
	  } 
	   } 
	   } 
	 & \\ 
 

    \theutterance \stepcounter{utterance}  

    & \multicolumn{4}{p{0.6\linewidth}}{\cellcolor[rgb]{0.8,1,0.9}{%
	\makecell[{{p{\linewidth}}}]{% 
	  \tt {\tiny [A$\rangle$GM]}  
	 RAGIONAMENTO STRATEGICO: {'Valuto il valore per me della nuova proposta: {A87(215), B33(30), A71(65), C17(138), B35(121), C00(262), C01(389), B86(500), A04(484)} = 2204. Sforzo: 4+594+175+313+635+461+539+399+203=3323. È molto inferiore alla mia proposta (3819) e alla variante di compromesso che avevo in mente (~3563). Devo spingere per reintrodurre i miei altovalore C54, A03 e possibilmente A10/B41/A77. Offro un compromesso forte senza B41: {A87,A04,A10,C54,B86,A03,A77} aveva sforzo 3361-720=2641? Ricalcolo: base full era 3361 con B41 incluso. Togliendo B41: 3361-720=2641. Aggiungo C01(539) => 3180 e ho ancora margine 206. Posso aggiungere A71(175) per 3355 (ok). Valore per me: 215+484+400+668+500+461+508+389+65=3690, molto alto ma potrebbe essere più digeribile per loro perché include C01 e A71 dalla loro proposta e non include B41. Verifico sforzo esatto del set compromesso: {A87, A04, A10, C54, B86, A03, A77, C01, A71} = 4+203+437+369+399+544+685+539+175=3355 ≤ 3386. Ottimo. Proporrò questo come compromesso. Argomenterò su equilibrio e inclusione di molte loro scelte.'}\\ \tt \\ \tt RIFIUTO: {'A87', 'B33', 'A71', 'C17', 'B35', 'C00', 'C01', 'B86', 'A04'}\\ \tt \\ \tt PROPOSTA: {'A87', 'A04', 'A10', 'C54', 'B86', 'A03', 'A77', 'C01', 'A71'}\\ \tt \\ \tt ARGOMENTAZIONE: {'Questa è una soluzione di compromesso che resta entro il limite (sforzo totale 3355 ≤ 3386) e incorpora diversi elementi della tua proposta (A87, A04, C01, A71, B86), affiancandoli a componenti ad alta efficienza complessiva (C54, A03, A10, A77). In questo modo bilanciamo pezzi “portanti” con moduli compatti, sfruttando al meglio il budget senza ricorrere a elementi eccessivamente pesanti. Se ritieni di dover ancora alleggerire, posso valutare la sostituzione di A77 con un blocco più leggero mantenendo il vincolo.'} 
	  } 
	   } 
	   } 
	 & & \\ 
 

    \theutterance \stepcounter{utterance}  

    & & \multicolumn{4}{p{0.6\linewidth}}{\cellcolor[rgb]{0.9,0.9,0.9}{%
	\makecell[{{p{\linewidth}}}]{% 
	  \tt {\tiny [GM$\rangle$B]}  
	 RIFIUTO: {'A87', 'B33', 'A71', 'C17', 'B35', 'C00', 'C01', 'B86', 'A04'}\\ \tt \\ \tt PROPOSTA: {'A87', 'A04', 'A10', 'C54', 'B86', 'A03', 'A77', 'C01', 'A71'}\\ \tt \\ \tt ARGOMENTAZIONE: {'Questa è una soluzione di compromesso che resta entro il limite (sforzo totale 3355 ≤ 3386) e incorpora diversi elementi della tua proposta (A87, A04, C01, A71, B86), affiancandoli a componenti ad alta efficienza complessiva (C54, A03, A10, A77). In questo modo bilanciamo pezzi “portanti” con moduli compatti, sfruttando al meglio il budget senza ricorrere a elementi eccessivamente pesanti. Se ritieni di dover ancora alleggerire, posso valutare la sostituzione di A77 con un blocco più leggero mantenendo il vincolo.'} 
	  } 
	   } 
	   } 
	 & \\ 
 

    \theutterance \stepcounter{utterance}  

    & & & \multicolumn{4}{p{0.6\linewidth}}{\cellcolor[rgb]{1,0.85,0.72}{%
	\makecell[{{p{\linewidth}}}]{% 
	  \tt {\tiny [GM$\langle$B]}  
	 RAGIONAMENTO STRATEGICO: {'Calcolo il mio valore sulla loro nuova proposta: {A87, A04, A10, C54, B86, A03, A77, C01, A71}. Valore per me = 461+138+262+30+121+215+97+389+583 = 2,296; sforzo 3355. È meglio della loro prima (1,389) ma ancora molto inferiore alla mia combinazione ottimale ~3,885. Voglio recuperare i miei pezzi chiave ad alto valore (B33, C17, B35, C00) sostituendo i loro elementi per me deboli (C54, A77, B41, A03 talvolta). Mantengo più possibile i loro preferiti: hanno insistito su A03, A10, A77, C54, B86, A71, C01, A04, A87. Posso proporre un compromesso forte per me e vicino ai loro: includo A87, A04, A10, A03, A71, C01 e aggiungo i miei B33, C17, C00. Verifico sforzo: A87(4)+A04(203)+A10(437)+A03(544)+A71(175)+C01(539)+B33(594)+C17(313)+C00(461)=3,270 entro 3,386. Valore per me: 461+138+262+215+583+389+668+484+400=3,600, ottimo. Lascio fuori A77, C54, B86 che per me valgono poco. Argomento che è entro limite, conserva gran parte dei loro elementi e massimizza coesione. Mi preparo ad offrire scambi alternativi se chiedono di reinserire uno tra B86 o C54: potrei valutare sostituzioni minori come scambiare A03 con B86 (sforzo -145, peggiora valore ma potrei chiedere qualcosa in cambio) o A04 con B86 (sforzo +196, totale 3,466 no), quindi meglio A03↔B86: nuovo sforzo 3,270-544+399=3,125 (accettabile) con valore che scende di 94. Primo passo: rifiuto e propongo il mio set 3,270 con apertura a uno scambio 1-a-1 su A03.'}\\ \tt \\ \tt RIFIUTO: {'A87', 'A04', 'A10', 'C54', 'B86', 'A03', 'A77', 'C01', 'A71'}\\ \tt \\ \tt PROPOSTA: {'A87', 'A04', 'A10', 'A03', 'A71', 'C01', 'B33', 'C17', 'C00'}\\ \tt \\ \tt ARGOMENTAZIONE: {'Questa selezione resta entro il limite (sforzo totale 3270 ≤ 3386) e mantiene la maggior parte degli elementi che hai proposto (A87, A04, A10, A03, A71, C01), affiancandoli a tre moduli strutturali che migliorano la robustezza complessiva (B33, C17, C00). Così evitiamo componenti meno incisivi come A77 o C54 pur restando molto compatti. Se per te è importante reinserire un elemento tra B86 o C54, posso valutare uno scambio 1‑a‑1 con A03, mantenendo il vincolo di sforzo.'} 
	  } 
	   } 
	   } 
	  \\ 
 

    \theutterance \stepcounter{utterance}  

    & & \multicolumn{4}{p{0.6\linewidth}}{\cellcolor[rgb]{0.9,0.9,0.9}{%
	\makecell[{{p{\linewidth}}}]{% 
	  \tt {\tiny [A$\langle$GM]}  
	 RIFIUTO: {'A87', 'A04', 'A10', 'C54', 'B86', 'A03', 'A77', 'C01', 'A71'}\\ \tt \\ \tt PROPOSTA: {'A87', 'A04', 'A10', 'A03', 'A71', 'C01', 'B33', 'C17', 'C00'}\\ \tt \\ \tt ARGOMENTAZIONE: {'Questa selezione resta entro il limite (sforzo totale 3270 ≤ 3386) e mantiene la maggior parte degli elementi che hai proposto (A87, A04, A10, A03, A71, C01), affiancandoli a tre moduli strutturali che migliorano la robustezza complessiva (B33, C17, C00). Così evitiamo componenti meno incisivi come A77 o C54 pur restando molto compatti. Se per te è importante reinserire un elemento tra B86 o C54, posso valutare uno scambio 1‑a‑1 con A03, mantenendo il vincolo di sforzo.'} 
	  } 
	   } 
	   } 
	 & \\ 
 

    \theutterance \stepcounter{utterance}  

    & \multicolumn{4}{p{0.6\linewidth}}{\cellcolor[rgb]{0.8,1,0.9}{%
	\makecell[{{p{\linewidth}}}]{% 
	  \tt {\tiny [A$\rangle$GM]}  
	 RAGIONAMENTO STRATEGICO: {'Valuto il valore per me della loro nuova proposta: {A87(215), A04(484), A10(400), A03(461), A71(65), C01(389), B33(30), C17(138), C00(262)} = 2444. Sforzo 3270. È ancora molto lontano dal mio obiettivo e omette due dei miei top (C54=668, B86=500) e anche A77/B41. Devo spingere per reinserire C54 e B86 almeno. Provo un contro-compromesso che includa molte delle loro scelte ma rimpiazzi i pezzi per me poco utili (B33, C17, C00, A71) con C54 e B86, restando entro il limite e mantenendo A87, A04, A10, A03, C01 che loro desiderano. Verifico set: {A87, A04, A10, A03, C01, C54, B86}. Sforzo: 4+203+437+544+539+369+399=2495. Ho molto margine: posso aggiungere A77 (685) => 3180; rimane 206. Posso anche aggiungere A71 (175) per venire incontro a loro => 3355 (ok). Valore per me di {A87, A04, A10, A03, C01, C54, B86, A77, A71} = 215+484+400+461+389+668+500+508+65 = 3690, molto alto per me, ma include molti dei loro (A87,A04,A10,A03,C01,A71) e aggiunge solo C54 e B86 e A77. Se respingono A77, ho una variante senza A77: {A87,A04,A10,A03,C01,C54,B86,A71} sforzo 3355-685=2670; posso proporla come opzione più leggera ma poi potrei chiedere di rimettere C17? Non mi serve. Meglio proporre direttamente una versione forte ma senza A77 per aumentarne l’accettabilità: {A87,A04,A10,A03,C01,C54,B86,A71} con sforzo 300? Ricalcolo esatto: 4+203+437+544+539+369+399+175=2670. Tanto margine; potrei aggiungere C17 (313) => 2983 mantenendo valore extra 138, ancora molto sotto il limite, ma non devo dare spazio a pezzi a basso valore per me se non necessario. Tuttavia lasciare tanto margine potrebbe sembrare sprecato; potrei inserire C00 (461) arrivando a 3131 e ancora sotto; il mio valore +262. Ma preferisco tenere la proposta compatta con altissima efficienza e accennare che c’è margine per aggiunte a loro scelta, purché non rimuovano C54 o B86. Quindi formulo proposta: {A87,A04,A10,A03,C01,C54,B86,A71}. Argomento che incorpora quasi tutta la loro lista tranne i tre moduli che per me sono deboli (B33,C17,C00), aggiungendo due ad altissima efficienza (C54,B86), e lasciando ampio margine che possiamo usare per un ultimo inserimento se vogliono, ma solo scambiando qualcosa di meno incisivo, o aggiungendo un pezzo leggero come C17 se proprio insistono.'}\\ \tt \\ \tt RIFIUTO: {'A87', 'A04', 'A10', 'A03', 'A71', 'C01', 'B33', 'C17', 'C00'}\\ \tt \\ \tt PROPOSTA: {'A87', 'A04', 'A10', 'A03', 'C01', 'C54', 'B86', 'A71'}\\ \tt \\ \tt ARGOMENTAZIONE: {'Rimaniamo entro il limite con ampio margine (sforzo totale 2670 ≤ 3386). Mantengo quasi tutti i componenti che proponi (A87, A04, A10, A03, C01, A71) e sostituisco i tre moduli meno incisivi (B33, C17, C00) con due elementi ad altissima efficienza complessiva (C54, B86). Questo aumenta l’impatto senza appesantire. Se desideri ancora aggiungere qualcosa, possiamo valutare un singolo inserimento extra rimanendo nel vincolo, ma senza rimuovere C54 o B86.'} 
	  } 
	   } 
	   } 
	 & & \\ 
 

    \theutterance \stepcounter{utterance}  

    & & \multicolumn{4}{p{0.6\linewidth}}{\cellcolor[rgb]{0.9,0.9,0.9}{%
	\makecell[{{p{\linewidth}}}]{% 
	  \tt {\tiny [GM$\rangle$B]}  
	 RIFIUTO: {'A87', 'A04', 'A10', 'A03', 'A71', 'C01', 'B33', 'C17', 'C00'}\\ \tt \\ \tt PROPOSTA: {'A87', 'A04', 'A10', 'A03', 'C01', 'C54', 'B86', 'A71'}\\ \tt \\ \tt ARGOMENTAZIONE: {'Rimaniamo entro il limite con ampio margine (sforzo totale 2670 ≤ 3386). Mantengo quasi tutti i componenti che proponi (A87, A04, A10, A03, C01, A71) e sostituisco i tre moduli meno incisivi (B33, C17, C00) con due elementi ad altissima efficienza complessiva (C54, B86). Questo aumenta l’impatto senza appesantire. Se desideri ancora aggiungere qualcosa, possiamo valutare un singolo inserimento extra rimanendo nel vincolo, ma senza rimuovere C54 o B86.'} 
	  } 
	   } 
	   } 
	 & \\ 
 

    \theutterance \stepcounter{utterance}  

    & & & \multicolumn{4}{p{0.6\linewidth}}{\cellcolor[rgb]{1,0.85,0.72}{%
	\makecell[{{p{\linewidth}}}]{% 
	  \tt {\tiny [GM$\langle$B]}  
	 RAGIONAMENTO STRATEGICO: {'La loro nuova proposta include 8 oggetti: {A87, A04, A10, A03, C01, C54, B86, A71} con sforzo 2670. Il mio valore per questo set è 461+138+262+215+389+30+121+583=2,199. Abbiamo ancora 716 di margine, quindi posso provare a inserire uno o due dei miei pezzi chiave ad alto valore (B33 594, C17 313, C00 461, B35 635, B32 695). Loro hanno chiesto esplicitamente di non rimuovere C54 o B86; posso accettarlo per ora. Mi conviene proporre un’aggiunta netta rimanendo entro il limite, senza togliere i loro due elementi preferiti. Con 716 posso aggiungere B33 (594) e restare a 3264, oppure C17 (313) e restare a 2983; potrei anche aggiungere entrambi? 2670+594+313=3577, supera il limite. Potrei aggiungere B33 da solo: miglior valore per me +668. Oppure aggiungere C17 da solo: +484. Qual è il massimo valore aggiungibile entro 716? B35 (635) darebbe +500; B33 è migliore (+668). Posso proporre di aggiungere B33 mantenendo tutto il resto invariato: sforzo totale 3264 ≤ 3386. Questo mi porta molto vicino ai miei obiettivi senza toccare C54 o B86. Argomento che sfruttiamo meglio il budget residuo. Se rifiutano, seconda opzione: aggiungere C17. Parto offrendo l’aggiunta di B33. Se puntano a un ulteriore “singolo inserimento extra”, posso anche accettare di aggiungere C17 al posto di A03 (ma non necessario ora). Presento proposta con B33 incluso.'}\\ \tt \\ \tt RIFIUTO: {'A87', 'A04', 'A10', 'A03', 'C01', 'C54', 'B86', 'A71'}\\ \tt \\ \tt PROPOSTA: {'A87', 'A04', 'A10', 'A03', 'C01', 'C54', 'B86', 'A71', 'B33'}\\ \tt \\ \tt ARGOMENTAZIONE: {'Restiamo pienamente entro il limite sfruttando meglio il margine inutilizzato: sforzo totale 3264 ≤ 3386. Mantengo integralmente C54 e B86 come richiesto e aggiungo un unico modulo aggiuntivo per consolidare il set senza appesantire. Se preferisci un’alternativa equivalente in termini di carico, posso considerare l’inserimento di C17 al posto di B33, restando comunque entro il vincolo.'} 
	  } 
	   } 
	   } 
	  \\ 
 

    \theutterance \stepcounter{utterance}  

    & & \multicolumn{4}{p{0.6\linewidth}}{\cellcolor[rgb]{0.9,0.9,0.9}{%
	\makecell[{{p{\linewidth}}}]{% 
	  \tt {\tiny [A$\langle$GM]}  
	 RIFIUTO: {'A87', 'A04', 'A10', 'A03', 'C01', 'C54', 'B86', 'A71'}\\ \tt \\ \tt PROPOSTA: {'A87', 'A04', 'A10', 'A03', 'C01', 'C54', 'B86', 'A71', 'B33'}\\ \tt \\ \tt ARGOMENTAZIONE: {'Restiamo pienamente entro il limite sfruttando meglio il margine inutilizzato: sforzo totale 3264 ≤ 3386. Mantengo integralmente C54 e B86 come richiesto e aggiungo un unico modulo aggiuntivo per consolidare il set senza appesantire. Se preferisci un’alternativa equivalente in termini di carico, posso considerare l’inserimento di C17 al posto di B33, restando comunque entro il vincolo.'} 
	  } 
	   } 
	   } 
	 & \\ 
 

    \theutterance \stepcounter{utterance}  

    & \multicolumn{4}{p{0.6\linewidth}}{\cellcolor[rgb]{0.8,1,0.9}{%
	\makecell[{{p{\linewidth}}}]{% 
	  \tt {\tiny [A$\rangle$GM]}  
	 RAGIONAMENTO STRATEGICO: {'Valuto il valore per me della loro nuova proposta: {A87(215), A04(484), A10(400), A03(461), C01(389), C54(668), B86(500), A71(65), B33(30)} = 3212. Sforzo: 3264. È molto meglio delle loro proposte precedenti e include i miei due top C54 e B86 insieme ai miei preferiti A04, A03, A10 e C01. L’unico elemento poco utile per me è B33 (30/594). Posso cercare un piccolo miglioramento sostituendo B33 con C17 (138/313), come loro stessi hanno suggerito come opzione: lo sforzo passerebbe da 3264 -594 +313 = 2983, ampiamente entro il limite, e il mio valore salirebbe a 3320 (+108). Potrei anche suggerire di aggiungere un ulteriore pezzo leggero dato l’ampio margine, ma rischierei di riaprire la trattativa. Dato che la loro proposta attuale è già molto favorevole per me e contiene quasi tutto ciò che voglio, potrei accettarla immediatamente per chiudere con un ottimo esito. Tuttavia, provo una controproposta minima e ragionevole: stessa lista ma sostituire B33 con C17, appoggiandomi alla loro stessa alternativa indicata. Se rifiutano, sono pronto ad accettare la loro proposta attuale.'}\\ \tt \\ \tt RIFIUTO: {'A87', 'A04', 'A10', 'A03', 'C01', 'C54', 'B86', 'A71', 'B33'}\\ \tt \\ \tt PROPOSTA: {'A87', 'A04', 'A10', 'A03', 'C01', 'C54', 'B86', 'A71', 'C17'}\\ \tt \\ \tt ARGOMENTAZIONE: {'Accolgo la tua apertura: manteniamo integralmente C54 e B86 e ottimizziamo il margine sostituendo B33 con C17, restando ampiamente entro il vincolo (sforzo totale 2983 ≤ 3386). È una variante leggermente più efficiente senza appesantire il carico complessivo. Se per te è equivalente, chiudiamo così; in alternativa posso confermare la tua versione con B33.'} 
	  } 
	   } 
	   } 
	 & & \\ 
 

    \theutterance \stepcounter{utterance}  

    & & \multicolumn{4}{p{0.6\linewidth}}{\cellcolor[rgb]{0.9,0.9,0.9}{%
	\makecell[{{p{\linewidth}}}]{% 
	  \tt {\tiny [GM$\rangle$B]}  
	 RIFIUTO: {'A87', 'A04', 'A10', 'A03', 'C01', 'C54', 'B86', 'A71', 'B33'}\\ \tt \\ \tt PROPOSTA: {'A87', 'A04', 'A10', 'A03', 'C01', 'C54', 'B86', 'A71', 'C17'}\\ \tt \\ \tt ARGOMENTAZIONE: {'Accolgo la tua apertura: manteniamo integralmente C54 e B86 e ottimizziamo il margine sostituendo B33 con C17, restando ampiamente entro il vincolo (sforzo totale 2983 ≤ 3386). È una variante leggermente più efficiente senza appesantire il carico complessivo. Se per te è equivalente, chiudiamo così; in alternativa posso confermare la tua versione con B33.'} 
	  } 
	   } 
	   } 
	 & \\ 
 

    \theutterance \stepcounter{utterance}  

    & & & \multicolumn{4}{p{0.6\linewidth}}{\cellcolor[rgb]{1,0.85,0.72}{%
	\makecell[{{p{\linewidth}}}]{% 
	  \tt {\tiny [GM$\langle$B]}  
	 RAGIONAMENTO STRATEGICO: {'Valuto la loro ultima proposta: {A87, A04, A10, A03, C01, C54, B86, A71, C17}. Sforzo 2983 ≤ 3386. Valore per me = A87(461)+A04(138)+A10(262)+A03(215)+C01(389)+C54(30)+B86(121)+A71(583)+C17(484) = 2,683. C’è ancora margine di 403 sforzo: potrei tentare di aggiungere un mio pezzo chiave. Migliori opzioni: C00(461) non entra; B33(594) non entra; B35(635) non entra; B32(695) non entra. L’unico ad alto valore che entra è A77(685) no; B41(720) no; A04 già incluso; C54 già incluso. Quindi non posso migliorare con un’aggiunta singola. Posso proporre uno scambio per inserire C00 (461 di sforzo) rimuovendo qualcosa di basso valore per me come C54(369) o B86(399) ma 403+369=772 o 403+399=802 è margine, però non posso “aggiungere” senza scambio; dovrei chiedere modifica. Tuttavia hanno chiarito che non vogliono rimuovere C54 o B86 in precedenza; ora li hanno mantenuti e stanno cercando chiusura. Il set mi dà 2,683, significativamente migliore delle loro opzioni precedenti e temo che spingere ulteriormente possa farli tornare indietro. Accetto per chiudere con buon risultato.'}\\ \tt \\ \tt ACCORDO: {'A87', 'A04', 'A10', 'A03', 'C01', 'C54', 'B86', 'A71', 'C17'}\\ \tt \\ \tt ARGOMENTAZIONE: {'Accetto questa soluzione: rispetta pienamente il limite (2983 ≤ 3386) e sfrutta bene il margine disponibile mantenendo i pilastri comuni. Chiudiamo su questo set.'} 
	  } 
	   } 
	   } 
	  \\ 
 

    \theutterance \stepcounter{utterance}  

    & & & \multicolumn{2}{p{0.3\linewidth}}{\cellcolor[rgb]{0.95,0.95,0.95}{%
	\makecell[{{p{\linewidth}}}]{% 
	  \tt {\tiny [GM$|$GM]}  
	 game successful 
	  } 
	   } 
	   } 
	 & & \\ 
 

    \theutterance \stepcounter{utterance}  

    & & & \multicolumn{2}{p{0.3\linewidth}}{\cellcolor[rgb]{0.95,0.95,0.95}{%
	\makecell[{{p{\linewidth}}}]{% 
	  \tt {\tiny [GM$|$GM]}  
	 end game 
	  } 
	   } 
	   } 
	 & & \\ 
 

\end{supertabular}
}

\end{document}
